% Regression: t2_seljbraid — imported current-behavior fixture from the handoff corpus.
% Target: RMP 2011.12127 fig3_fig3-pepe_seljbraidanyon.pdf (topological
% spin / self-braiding of an anyonic MPO string).
% Formula:  ring_j wound twice = e^{i phi_j} ring_j  —  the anyon-j MPO
% ring wound TWICE around the site (drawn flat: two nested loops joined
% by a crossover near the marked tensor j, free end exiting) equals the
% phase e^{i phi_j} times the singly-wound ring.
% Genuinely non-grid (nested loops threading glyphs): tenkzfree; the
% double winding is stated topologically as outer loop -> crossover ->
% inner loop -> tail, which is the paper's own flattening.
% Honest gaps: \tnjoin has no wire-hue key (the ring is red in the
% paper), \tnput has no marked/hue skin for the blue j, the mpo skin is
% a house box, and the crossover renders as a plain transversal
% intersection (no over/under gap primitive).
\documentclass[varwidth=19cm, border=6pt]{standalone}
\usepackage{amsmath, amssymb}
\usepackage{tenkz}
\begin{document}

\[
  \begin{tenkzfree}
    % outer winding: four MPO tensors
    \tnput[mpo]{Wo}{(-24mm,0)}{}
    \tnput[mpo]{No}{(0,13mm)}{}
    \tnput[mpo]{Eo}{(24mm,0)}{}
    \tnput[mpo]{So}{(0,-13mm)}{}
    % inner winding: three MPO tensors and the marked tensor j
    \tnput[mpo]{Wi}{(-14mm,0)}{}
    \tnput[mpo]{Ni}{(0,7mm)}{}
    \tnput[mpo]{Ei}{(14mm,0)}{}
    \tnput[mpo]{J}{(6mm,-7mm)}{j}
    % ONE strand, wound twice; every MPO tensor carries exactly two
    % string legs.  Inner winding: J -> Wi -> Ni -> Ei ...
    \tnjoin[route=curve, out=180, in=-90]{J.west}{Wi.south}
    \tnjoin[route=curve, out=90, in=180]{Wi.north}{Ni.west}
    \tnjoin[route=curve, out=0, in=90]{Ni.east}{Ei.north}
    % ... crossover down into the outer winding (SE sector) ...
    \tnjoin[route=curve, out=-90, in=0]{Ei.south}{So.east}
    % ... outer winding: So -> Wo -> No -> Eo ...
    \tnjoin[route=curve, out=180, in=-90]{So.west}{Wo.south}
    \tnjoin[route=curve, out=90, in=180]{Wo.north}{No.west}
    \tnjoin[route=curve, out=0, in=90]{No.east}{Eo.north}
    % ... crossover back into J: closes the doubly-wound ring, crossing
    % the first SE connector once (the flattened self-braiding)
    \tnjoin[route=curve, out=-90, in=0]{Eo.south}{J.east}
    % the free string end at the marked tensor, crossing the outer loop
    \tnjoin[route=curve, out=-60, in=160]{J.south}{18mm,-19mm}
  \end{tenkzfree}
  \;=\; e^{i\phi_j}\;
  \begin{tenkzfree}
    % single winding: three MPO tensors and the marked tensor j
    \tnput[mpo]{W}{(-16mm,0)}{}
    \tnput[mpo]{N}{(0,9mm)}{}
    \tnput[mpo]{E}{(16mm,0)}{}
    \tnput[mpo]{J}{(4mm,-9mm)}{j}
    \tnjoin[route=curve, out=90, in=180]{W.north}{N.west}
    \tnjoin[route=curve, out=0, in=90]{N.east}{E.north}
    \tnjoin[route=curve, out=-90, in=0]{E.south}{J.east}
    \tnjoin[route=curve, out=180, in=-90]{J.west}{W.south}
    % free string end
    \tnjoin{J.south east}{12mm,-15mm}
  \end{tenkzfree}
\]

\end{document}
